% ============================================================================
% WORKSHEET - Section 1.7: Proportional Reasoning with Scale Models
% Oklahoma additional topic
% Included States: Oklahoma
% ============================================================================
\section{Proportional Reasoning with Scale Models}
\setcounter{wsNumber}{0}

\begin{worksheetSkills}
\begin{itemize}
    \item Use a scale factor to move between a model and a real object
    \item Convert measurements so the scale comparison uses matching units
    \item Interpret scale models in architecture, engineering, and design contexts
\end{itemize}
\end{worksheetSkills}

\worksheetHeader{Scale Model Sprint}

\begin{speedDrill}{Find the Real or Model Measure}
\timerBadge{4}

\renewcommand{\arraystretch}{1.9}
\begin{center}
\begin{tabular}{|C{7cm}|C{4cm}|}
\hline
\textbf{Problem} & \textbf{Answer} \\
\hline
Model car: $8$ cm, scale $1 : 25$; find actual length & \answerBlank[3cm] \\
\hline
Bridge: $450$ m actual, scale $1 : 500$; find model length & \answerBlank[3cm] \\
\hline
Room: $6$ m actual, scale $1 : 50$; find model length & \answerBlank[3cm] \\
\hline
Model tower: $12$ cm, scale $1 : 200$; find actual height & \answerBlank[3cm] \\
\hline
Real wing: $12$ m, model wing: $15$ cm; find the scale & \answerBlank[3cm] \\
\hline
\end{tabular}
\end{center}
\end{speedDrill}
\worksheetAnswer{(1)~Actual length $= 8 \times 25 = 200$ cm. (2)~Convert $450$ m to centimeters: $450$ m $= 45{,}000$ cm. Model length $= 45{,}000 \div 500 = 90$ cm. (3)~$6$ m $= 600$ cm, so the model length is $600 \div 50 = 12$ cm. (4)~Actual height $= 12 \times 200 = 2400$ cm $= 24$ m. (5)~Convert $12$ m to $1200$ cm. Then $1200 \div 15 = 80$, so the scale is $1 : 80$.}

\worksheetHeader{Blueprint Decisions}

\begin{mathScenario}{Architect's Sketchbook}

An architect uses a blueprint scale of $1 : 100$.

\bigskip
\scenarioItem{Door on blueprint}{$0.9$ cm wide}
\scenarioItem{Window on blueprint}{$1.5$ cm wide}
\scenarioItem{Hallway in real life}{$4$ m long}

\bigskip

How wide is the real door? \answerBlank[4cm]

\bigskip

How wide is the real window? \answerBlank[4cm]

\bigskip

How long should the hallway be on the blueprint? \answerBlank[4cm]
\end{mathScenario}
\worksheetAnswer{(1)~With a $1 : 100$ scale, multiply the blueprint measure by $100$: $0.9 \times 100 = 90$ cm, so the real door is $90$ cm wide. (2)~The real window is $1.5 \times 100 = 150$ cm wide. (3)~Convert $4$ m to $400$ cm, then divide by $100$: $400 \div 100 = 4$ cm. The hallway should be $4$ cm long on the blueprint.}

\worksheetHeader{Scale Factor Mistakes}

\begin{errorDetective}{Multiply or Divide?}

\bigskip

Problem: A model bridge is $6$ cm long with scale $1 : 200$. Find the real length.
\errorWork{$6 \div 200 = 0.03$ cm}
Correct or wrong? \answerBlank[2cm] \quad Correct answer: \answerBlank[5cm]

\bigskip

Problem: A real road is $12$ m long with scale $1 : 300$. Find the model length.
\errorWork{$12 \times 300 = 3600$ m}
Correct or wrong? \answerBlank[2cm] \quad Correct answer: \answerBlank[5cm]

\bigskip

Problem: A model airplane wing is $20$ cm and the real wing is $10$ m.
\errorWork{Scale $= 20 : 10 = 2 : 1$}
Correct or wrong? \answerBlank[2cm] \quad Correct scale: \answerBlank[5cm]
\end{errorDetective}
\worksheetAnswer{All three are wrong. (1)~To go from model to real, multiply by the scale factor: $6 \times 200 = 1200$ cm, so the real length is $12$ m. (2)~To go from real to model, first convert to the same units: $12$ m $= 1200$ cm. Then divide by $300$: $1200 \div 300 = 4$ cm. (3)~The units must match. Convert $10$ m to $1000$ cm. Then compare model to real: $20 : 1000 = 1 : 50$.}

\worksheetHeader{Explain the Unit Conversion}

\begin{explainIt}{Why Must the Units Match?}

Explain why scale problems can go wrong if one measurement is in meters and the other is in centimeters.

\writeLines{5}

\bigskip

Describe how you decide whether to multiply or divide by the scale factor when solving a scale-model problem.

\writeLines{6}
\end{explainIt}
\worksheetAnswer{(1)~The units must match because a scale compares one measurement directly to another. If one value is in meters and the other is in centimeters, the numbers are not describing the same size units. For example, $10$ meters is not just $10$ centimeters; it is $1000$ centimeters. Without converting, the scale will be completely wrong. (2)~Multiply by the scale factor when moving from a small model to the larger real object. Divide by the scale factor when moving from the real object back to the smaller model. A good question is: am I scaling up or scaling down?}

\worksheetDone